\documentclass[11pt,a4paper]{article}
\usepackage[margin=1in]{geometry}
\usepackage{times}
\usepackage{latexsym}
\usepackage{amsmath}
\usepackage{amssymb}
\usepackage{booktabs}
\usepackage{multirow}
\usepackage{graphicx}
\usepackage{array}
\usepackage{xcolor}
\usepackage[colorlinks=true,citecolor=blue,urlcolor=blue]{hyperref}

\title{Output-Distribution Distillation Makes Compensation Networks Obsolete\\for SwiGLU Activation Sparsity}

\author{Anonymous}

\begin{document}
\maketitle

% ============================================================
\begin{abstract}
Activation sparsity in SwiGLU-based large language models enables significant inference speedups by dynamically pruning neurons, but existing learned approaches (e.g., FastForward) train per-neuron binary predictors with BCE loss and require auxiliary compensation networks ($\sim$256\,MB) to recover quality---adding latency and memory overhead.
We show that replacing BCE with KL-divergence distillation on the output distribution \emph{eliminates the need for compensation entirely}.
On Llama~3.1~8B at 50\% model-level sparsity, a KL-trained contextual predictor without compensation achieves 9.20 perplexity on WikiText-2 (within 50\% of the dense 6.12 baseline), while BCE with compensation reaches only 18.64---a gap of 9.44 PPL points.
The advantage stems from KL's gradient structure, which couples all active neurons through the output Jacobian, capturing joint interaction effects that BCE's factorized per-neuron gradient misses.
We formalize a modular decomposition $\mathbf{c} = \mathbf{c}_s + \mathbf{c}_\theta(\mathbf{x})$ combining SPON's static bias correction with KL contextual prediction, achieving 7.98 PPL with only 6.8\% interaction residual.
A seven-level sparsity sweep confirms the theoretically predicted quadratic scaling of the KL--BCE quality gap ($R^2 = 0.97$).
Cross-architecture validation on Mistral-7B-v0.3, comprehensive downstream benchmarks (MMLU, GSM8K, HumanEval, LongBench, MT-Bench, RULER), and latency profiling demonstrate that KL-only Pareto-dominates BCE+compensation: better quality, lower latency (12.0 vs.\ 13.8\,ms/token), and 10$\times$ fewer additional parameters (24\,MB vs.\ 280\,MB).
\end{abstract}

% ============================================================
\section{Introduction}

SwiGLU feed-forward networks~\cite{shazeer2020glu} in modern LLMs exhibit substantial contextual activation sparsity: at any given input, a large fraction of intermediate neurons contribute negligibly to the output.
Exploiting this sparsity at inference time by dynamically skipping inactive neurons can yield wall-clock speedups of $1.3$--$1.5\times$, a practically significant reduction for latency-sensitive deployment.

Three tiers of approaches have emerged.
\emph{Training-free} methods (TEAL, WINA, R-Sparse) apply static heuristics---magnitude thresholds, weight-informed scores---to select neurons.
While requiring no training, these methods incur a 6--9 PPL point gap at 50\% sparsity on Llama~3.1~8B.
\emph{Static learned} approaches such as SPON learn input-independent per-layer bias corrections via KL distillation, recovering average-case activation drift but unable to adapt to per-input variation.
\emph{Contextual learned} methods, exemplified by FastForward, train a lightweight MLP to predict per-token neuron importance.
However, FastForward uses binary cross-entropy (BCE) to train this predictor independently per neuron, then attaches a compensation network to correct the residual error---adding $\sim$256\,MB of parameters and a per-layer GEMV at inference time.

We ask: \emph{does the compensation network compensate for a fundamental limitation of the predictor, or for a limitation of the training objective?}
The answer hinges on gradient structure.
The BCE gradient $\partial\mathcal{L}_{\mathrm{BCE}}/\partial s_j = (\sigma(s_j) - t_j)\sigma'(s_j)$ depends only on neuron $j$'s own target, factorizing across neurons.
In contrast, the KL gradient $\partial\mathcal{L}_{\mathrm{KL}}/\partial s_j$ propagates through the full Jacobian $\partial\mathbf{z}/\partial\mathbf{h} \cdot \partial\mathbf{h}/\partial\mathbf{m}$, coupling all active neurons.
This coupling captures joint interaction effects that scale as $O(s^2 D^2)$ with sparsity level $s$ and hidden dimension $D$---precisely the effects that compensation networks attempt to correct post hoc.

Our contributions:
\begin{enumerate}
    \item We demonstrate that KL-trained contextual sparsity predictors categorically dominate BCE+compensation: 9.20 vs.\ 18.64 PPL on Llama~3.1~8B at 50\% sparsity, with lower latency and 10$\times$ fewer parameters (\S\ref{sec:ablation}, \S\ref{sec:fullbench}).
    \item We formalize a modular decomposition $\mathbf{c} = \mathbf{c}_s + \mathbf{c}_\theta(\mathbf{x})$ combining static (SPON) and contextual (KL) corrections, achieving 7.98 PPL with near-additive composition (6.8\% interaction residual; \S\ref{sec:composition}).
    \item We provide a testable theoretical prediction---the KL--BCE quality gap scales quadratically with sparsity---and confirm it empirically ($R^2 = 0.97$ across seven sparsity levels; \S\ref{sec:scaling}).
\end{enumerate}

% ============================================================
\section{Related Work}

\paragraph{Training-free activation sparsity.}
TEAL~\cite{liu2024teal} uses magnitude-based thresholding with greedy per-layer budget allocation.
WINA~\cite{shen2024wina} incorporates weight magnitudes into the scoring criterion.
R-Sparse~\cite{zhang2024rsparse} introduces rank-aware scoring.
All three operate without any parameter training but suffer quality gaps of 6--9 PPL points at 50\% sparsity.

\paragraph{Learned sparsity prediction.}
FastForward~\cite{wang2024fastforward} trains a small MLP per layer to predict binary neuron masks using BCE, then adds a compensation MLP to correct residual errors.
SPON~\cite{zhou2024spon} takes a different approach: learning input-independent per-layer bias vectors via KL distillation of the teacher's output distribution.
ShadowLLM~\cite{chen2024shadowllm} uses MSE on intermediate representations for per-layer correction.

\paragraph{Knowledge distillation for sparsity.}
Output-distribution matching via KL divergence is standard in knowledge distillation~\cite{hinton2015distilling} but has not been applied to train \emph{contextual} sparsity predictors.
The key distinction is that we distill the \emph{sparsity decision} (which neurons to activate) rather than the model weights, and the KL gradient structure provides a principled advantage over per-neuron objectives.

% ============================================================
\section{Method}
\label{sec:method}

\subsection{Problem Setup}

Consider a SwiGLU FFN layer with gate projection $\mathbf{g} = \sigma_{\mathrm{SiLU}}(\mathbf{W}_g \mathbf{x})$, up projection $\mathbf{u} = \mathbf{W}_u \mathbf{x}$, and element-wise product $\mathbf{h} = \mathbf{g} \odot \mathbf{u}$.
A binary mask $\mathbf{m} \in \{0,1\}^D$ selects a subset of neurons, yielding sparse output $\mathbf{y}_{\mathrm{sparse}} = \mathbf{W}_d (\mathbf{m} \odot \mathbf{h})$.
We seek a predictor $f_\theta(\mathbf{x}): \mathbb{R}^{d} \to [0,1]^D$ that produces soft scores, relaxed to masks via Gumbel-sigmoid sampling during training and hard thresholding at inference.

\subsection{Gradient Coupling: KL vs.\ BCE}

\paragraph{BCE objective.}
Given oracle mask $\mathbf{t}$ (e.g., top-$k$ by activation magnitude), the BCE loss is:
\begin{equation}
\mathcal{L}_{\mathrm{BCE}} = -\sum_{j=1}^{D} \big[ t_j \log \sigma(s_j) + (1-t_j)\log(1-\sigma(s_j)) \big]
\end{equation}
where $s_j = f_\theta(\mathbf{x})_j$.
The gradient $\partial\mathcal{L}_{\mathrm{BCE}}/\partial s_j = \sigma(s_j) - t_j$ depends only on neuron $j$---the loss \emph{factorizes} across neurons.

\paragraph{KL objective.}
We instead minimize the KL divergence between the dense and sparse model output distributions:
\begin{equation}
\mathcal{L}_{\mathrm{KL}} = \mathrm{KL}\!\left(p_{\mathrm{dense}}(\cdot|\mathbf{x}) \;\|\; p_{\mathrm{sparse}}(\cdot|\mathbf{x}, \mathbf{m})\right)
\end{equation}
The gradient with respect to score $s_j$ is:
\begin{equation}
\frac{\partial\mathcal{L}_{\mathrm{KL}}}{\partial s_j} = \underbrace{\frac{\partial\mathcal{L}_{\mathrm{KL}}}{\partial \mathbf{z}}}_{\text{output error}} \cdot \underbrace{\frac{\partial \mathbf{z}}{\partial \mathbf{h}}}_{\text{down proj}} \cdot \underbrace{\frac{\partial \mathbf{h}}{\partial m_j}}_{\text{neuron }j} \cdot \frac{\partial m_j}{\partial s_j}
\end{equation}
Critically, $\partial\mathbf{z}/\partial\mathbf{h} = \mathbf{W}_d$ and $\partial\mathcal{L}_{\mathrm{KL}}/\partial\mathbf{z}$ depend on the \emph{full} sparse output, which in turn depends on \emph{all} active neurons.
Thus the gradient for neuron $j$ is \emph{coupled} to all other neurons' mask decisions through the shared output space.

\paragraph{Quadratic scaling prediction.}
When sparsity is $s$, each pruned neuron interacts with $O(sD)$ other pruned neurons.
The total pairwise interaction effects scale as $O(s^2 D^2)$.
BCE's factorized gradient ignores all such interactions, while KL captures them through the Jacobian.
This yields the testable prediction that the KL--BCE quality gap $\Delta(s) \approx a s^2$ for some constant $a$.

\subsection{Modular Decomposition}
\label{sec:decomposition}

We decompose the total correction to the sparse hidden states as:
\begin{equation}
\mathbf{c} = \mathbf{c}_s + \mathbf{c}_\theta(\mathbf{x})
\end{equation}
where $\mathbf{c}_s = \mathbb{E}[\Delta\mathbf{h}]$ is the static, input-independent bias correction (as in SPON), and $\mathbf{c}_\theta(\mathbf{x}) = \Delta\mathbf{h}(\mathbf{x}) - \mathbb{E}[\Delta\mathbf{h}]$ is the per-input residual variation captured by the KL contextual predictor.
Since $\mathbf{c}_s$ is constant and $\mathbf{c}_\theta(\mathbf{x})$ is zero-mean by construction, the two corrections are approximately orthogonal.
We measure the interaction residual as:
\begin{equation}
\rho = \frac{\|\mathbf{c} - (\mathbf{c}_s + \mathbf{c}_\theta(\mathbf{x}))\|_2}{\|\mathbf{c}_s\|_2 + \|\mathbf{c}_\theta(\mathbf{x})\|_2}
\end{equation}
and predict $\rho \ll 1$ when SPON and the KL predictor address different error components.

% ============================================================
\section{Experiments}
\label{sec:experiments}

\subsection{Setup}

\paragraph{Models.}
We evaluate on Llama~3.1~8B (primary) and Mistral-7B-v0.3 (cross-architecture validation).
Both use SwiGLU FFN layers with intermediate dimension $D = 14336$ (Llama) and $D = 14336$ (Mistral).

\paragraph{Sparsity target.}
Unless otherwise noted, we target 50\% model-level sparsity with per-layer budget allocation following TEAL's greedy algorithm.

\paragraph{Training.}
KL and BCE predictors share identical architecture (2-layer MLP, $\sim$24\,MB per model).
Training uses 1000 steps on a 10M-token calibration subset of C4.
Gumbel-sigmoid relaxation ($\tau = 1.0 \to 0.1$ cosine anneal) enables gradient flow through discrete masks during training; hard thresholds are used at inference.
We report WikiText-2 perplexity as the primary metric.

\paragraph{Baselines.}
We compare against: (1)~Dense (no sparsity), (2)~TEAL, WINA, R-Sparse (training-free), (3)~SPON (static learned), (4)~BCE-only, BCE+compensation (contextual learned, FastForward-style), and (5)~our KL-only, KL+compensation, and SPON+KL variants.

% ------------------------------------------------------------
\subsection{Four-Condition Ablation}
\label{sec:ablation}

Table~\ref{tab:ablation} presents the core ablation at 50\% sparsity on Llama~3.1~8B.

\begin{table}[t]
\centering
\small
\caption{Four-condition ablation on WikiText-2 (Llama~3.1~8B, 50\% sparsity). KL-only outperforms BCE+compensation by 9.44 PPL points without any compensation network.}
\label{tab:ablation}
\begin{tabular}{lcccc}
\toprule
\textbf{Condition} & \textbf{Seeds} & \textbf{PPL} & \textbf{$\Delta$ Dense} & \textbf{Comp.\ Benefit} \\
\midrule
Dense & -- & 6.12 & -- & -- \\
BCE-only & 3 & 28.16\small{$\pm$0.17} & +360\% & -- \\
BCE+comp & 5 & 18.64\small{$\pm$0.33} & +205\% & $-$9.52 \\
KL-only & 5 & 9.20\small{$\pm$0.04} & +50\% & -- \\
KL+comp & 3 & 9.02\small{$\pm$0.07} & +47\% & $-$0.18 \\
\bottomrule
\end{tabular}
\end{table}

KL-only achieves 9.20$\pm$0.04 PPL, surpassing BCE+compensation (18.64$\pm$0.33) by 9.44 points.
Adding compensation to KL yields only 0.18 PPL improvement (9.02 vs.\ 9.20), compared to 9.52 for BCE.
The extremely low standard deviation across 5 seeds (0.04) confirms high reproducibility.
The compensation network is made redundant by KL's gradient coupling---it corrects at training time what compensation would correct at inference time.

% ------------------------------------------------------------
\subsection{Isolating Loss vs.\ Relaxation (STE Baseline)}
\label{sec:ste}

A natural concern is whether the advantage comes from the KL loss function or from Gumbel-sigmoid relaxation.
We replace Gumbel-sigmoid with Straight-Through Estimation (STE) while keeping KL loss, shown in Table~\ref{tab:ste}.

\begin{table}[t]
\centering
\small
\caption{Relaxation method ablation. STE+KL still dominates Gumbel+BCE, confirming the loss function drives $>$97\% of the advantage.}
\label{tab:ste}
\begin{tabular}{llcc}
\toprule
\textbf{Condition} & \textbf{Relaxation} & \textbf{Loss} & \textbf{PPL} \\
\midrule
Gumbel+KL (ours) & Gumbel-sigmoid & KL & 9.20\small{$\pm$0.04} \\
STE+KL & Straight-Through & KL & 9.74\small{$\pm$0.09} \\
Gumbel+BCE & Gumbel-sigmoid & BCE & 28.16\small{$\pm$0.17} \\
\bottomrule
\end{tabular}
\end{table}

STE+KL (9.74) is slightly worse than Gumbel+KL (9.20, $\Delta = 0.54$) but dramatically better than Gumbel+BCE (28.16).
The Gumbel relaxation accounts for only 2.8\% of the total KL advantage over BCE (0.54 out of 18.96 PPL points), confirming that the gradient coupling property of KL---not the relaxation method---drives the quality improvement.

% ------------------------------------------------------------
\subsection{Iso-Compute Comparison}
\label{sec:isocompute}

KL requires a full forward pass per step to compute logit targets ($\sim$3.2$\times$ BCE per-step FLOPs).
To control for this, we train BCE for 3200 steps (equal total FLOPs), shown in Table~\ref{tab:isocompute}.

\begin{table}[t]
\centering
\small
\caption{Iso-compute comparison. Extended BCE training (3.2$\times$ steps) barely closes the gap.}
\label{tab:isocompute}
\begin{tabular}{lccc}
\toprule
\textbf{Condition} & \textbf{Steps} & \textbf{Total FLOPs} & \textbf{PPL} \\
\midrule
KL (Gumbel) & 1000 & 3200$\times$ & 9.20\small{$\pm$0.04} \\
BCE (standard) & 1000 & 1000$\times$ & 28.16\small{$\pm$0.17} \\
BCE (extended) & 3200 & 3200$\times$ & 25.41\small{$\pm$0.31} \\
\bottomrule
\end{tabular}
\end{table}

Extended BCE training reduces PPL from 28.16 to 25.41 ($-$2.75), but the remaining gap to KL is 16.21 points.
BCE's information bottleneck is at the loss function level: more compute cannot overcome a factorized gradient that ignores joint neuron interactions.

% ------------------------------------------------------------
\subsection{Training-Free Baselines and Quality Hierarchy}
\label{sec:baselines}

Table~\ref{tab:hierarchy} establishes the three-tier quality hierarchy across all method categories.

\begin{table}[t]
\centering
\small
\caption{Quality hierarchy at 50\% sparsity (Llama~3.1~8B, WikiText-2).}
\label{tab:hierarchy}
\begin{tabular}{llcc}
\toprule
\textbf{Method} & \textbf{Type} & \textbf{PPL} & \textbf{$\Delta$ Dense} \\
\midrule
Dense & -- & 6.12 & -- \\
TEAL & Training-free & 12.48 & +104\% \\
R-Sparse & Training-free & 13.62 & +123\% \\
WINA & Training-free & 14.76 & +141\% \\
SPON & Static learned & 10.82 & +77\% \\
KL-only & Contextual learned & 9.20 & +50\% \\
\bottomrule
\end{tabular}
\end{table}

The hierarchy is clean: training-free methods (12.5--14.8 PPL) $<$ static learned (10.82) $<$ contextual learned with KL (9.20).
Each tier corresponds to increasing information quality: heuristic scoring $\to$ average-case calibration $\to$ per-input output-distribution optimization.
TEAL, the best training-free method, is still 3.28 PPL points worse than KL-only, justifying the overhead of learned prediction.

% ------------------------------------------------------------
\subsection{SPON Validation}
\label{sec:spon}

We reimplement SPON on Llama~3.1~8B to serve as the static correction component.
Table~\ref{tab:spon} validates correctness via QA benchmarks.

\begin{table}[t]
\centering
\small
\caption{SPON reimplementation validation (Llama~3.1~8B, 50\% sparsity).}
\label{tab:spon}
\begin{tabular}{lccc}
\toprule
\textbf{Benchmark} & \textbf{Dense} & \textbf{SPON} & $\boldsymbol{\Delta}$ \\
\midrule
ARC-c & 55.2 & 51.2 & $-$4.0 \\
ARC-e & 79.4 & 74.5 & $-$4.9 \\
BoolQ & 83.1 & 79.3 & $-$3.8 \\
PIQA & 81.3 & 78.1 & $-$3.2 \\
WinoGrande & 77.8 & 74.8 & $-$3.0 \\
\midrule
WikiText-2 PPL & 6.12 & 10.82 & +77\% \\
\bottomrule
\end{tabular}
\end{table}

SPON achieves 10.82 PPL with accuracy drops of 3--5 points, consistent with the original paper's reported ranges.
As a static correction, SPON recovers the expected activation drift $\mathbb{E}[\Delta\mathbf{h}]$ but cannot adapt to per-input variation.

% ------------------------------------------------------------
\subsection{SPON + KL Composition}
\label{sec:composition}

We test the modular decomposition $\mathbf{c} = \mathbf{c}_s + \mathbf{c}_\theta(\mathbf{x})$ by applying SPON bias correction jointly with the KL predictor.
Table~\ref{tab:composition} reports the additivity analysis.

\begin{table}[t]
\centering
\small
\caption{Composition additivity analysis. SPON and KL corrections compose near-additively (6.8\% interaction residual in logit space).}
\label{tab:composition}
\begin{tabular}{lcccc}
\toprule
\textbf{Component} & \textbf{PPL Recovery} & \textbf{Pred.\ Additive} & \textbf{Actual} & \textbf{Residual} \\
\midrule
SPON only & 1.68 & -- & -- & -- \\
KL only & 3.30 & -- & -- & -- \\
SPON+KL & -- & 4.98 & 4.52 & 9.2\% \\
\bottomrule
\end{tabular}
\end{table}

SPON alone recovers 1.68 PPL points (from the TEAL baseline), KL alone recovers 3.30, and the combination recovers 4.52---close to the predicted additive 4.98 (9.2\% shortfall in PPL-space, 6.8\% in logit $L_2$ norm).
The resulting SPON+KL achieves 7.98$\pm$0.10 PPL, only 30.4\% above dense.

Per-layer cosine similarity between SPON and KL corrections confirms approximate orthogonality: mean cosine 0.11, median 0.08, with only 4 of 32 layers exceeding 0.20.
This validates the decomposition: SPON corrects the constant drift component, and KL corrects the per-input residual, with minimal interference.

% ------------------------------------------------------------
\subsection{Full Benchmark Suite}
\label{sec:fullbench}

Table~\ref{tab:fullbench} presents comprehensive results across perplexity and downstream tasks.

\begin{table*}[t]
\centering
\small
\caption{Comprehensive benchmark results (Llama~3.1~8B, 50\% sparsity). KL-only averages only 1.3\% accuracy drop from dense across five tasks; SPON+KL achieves 0.7\%.}
\label{tab:fullbench}
\begin{tabular}{lcccccccc}
\toprule
\textbf{Method} & \textbf{Wiki$\downarrow$} & \textbf{C4$\downarrow$} & \textbf{ARC-c$\uparrow$} & \textbf{WinoG$\uparrow$} & \textbf{MMLU$\uparrow$} & \textbf{GSM8K$\uparrow$} & \textbf{HumanE$\uparrow$} & \textbf{Avg $\Delta\downarrow$} \\
\midrule
Dense & 6.12 & 7.85 & 55.2 & 77.8 & 65.8 & 53.8 & 32.9 & --- \\
TEAL & 12.48 & 15.12 & 48.1 & 72.4 & 58.2 & 44.2 & 26.2 & $-$7.8 \\
WINA & 14.76 & 17.88 & 46.3 & 70.8 & 56.4 & 41.8 & 24.4 & $-$9.9 \\
R-Sparse & 13.62 & 16.31 & 47.2 & 71.5 & 57.1 & 43.1 & 25.6 & $-$8.8 \\
SPON & 10.82 & 13.24 & 51.2 & 74.8 & 60.3 & 46.5 & 28.1 & $-$5.3 \\
BCE+comp & 18.64 & 22.15 & 44.8 & 68.3 & 54.9 & 39.7 & 23.8 & $-$12.8 \\
KL-only & 9.20 & 11.38 & 54.3 & 76.9 & 65.0 & 52.8 & 32.3 & $-$1.3 \\
SPON+KL & 7.98 & 9.81 & 54.8 & 77.2 & 65.4 & 53.2 & 32.5 & $-$0.7 \\
\bottomrule
\end{tabular}
\end{table*}

KL-only's low perplexity (9.20 WikiText-2, 11.38 C4) transfers directly to downstream accuracy: the average drop across five tasks is only 1.3\%, compared to 12.8\% for BCE+compensation.
Tasks requiring precise multi-step reasoning (GSM8K, HumanEval) show the largest relative advantage for KL---BCE+compensation drops 26--28\% from dense on these tasks, while KL-only drops less than 2\%.
SPON+KL further closes the gap to 0.7\% average drop.

\paragraph{Inference efficiency.}
Table~\ref{tab:latency} presents the latency--quality Pareto analysis.

\begin{table}[t]
\centering
\small
\caption{Inference efficiency comparison. KL-only Pareto-dominates BCE+comp on all axes.}
\label{tab:latency}
\begin{tabular}{lcccc}
\toprule
\textbf{Method} & \textbf{ms/tok} & \textbf{Speedup} & \textbf{Extra MB} & \textbf{PPL} \\
\midrule
Dense & 16.8 & 1.00$\times$ & 0 & 6.12 \\
TEAL & 11.2 & 1.50$\times$ & 0 & 12.48 \\
SPON & 10.9 & 1.54$\times$ & 1 & 10.82 \\
KL-only & 12.0 & 1.40$\times$ & 24 & 9.20 \\
BCE+comp & 13.8 & 1.22$\times$ & 280 & 18.64 \\
SPON+KL & 12.1 & 1.39$\times$ & 25 & 7.98 \\
\bottomrule
\end{tabular}
\end{table}

KL-only strictly Pareto-dominates BCE+compensation: better quality (9.20 vs.\ 18.64 PPL), lower latency (12.0 vs.\ 13.8\,ms/token), and 10$\times$ fewer additional parameters (24 vs.\ 280\,MB).
The 1.8\,ms/token savings come from eliminating the per-layer compensation GEMV.
SPON+KL adds only 0.1\,ms (bias addition is nearly free) while further improving quality to 7.98 PPL.

% ------------------------------------------------------------
\subsection{Quadratic Scaling of KL--BCE Gap}
\label{sec:scaling}

We train KL-only and BCE-only at seven sparsity levels (30--60\%) with single seeds.
Table~\ref{tab:sweep} reports the results and fit statistics.

\begin{table}[t]
\centering
\small
\caption{Sparsity sweep: KL--BCE quality gap $\Delta$ follows $as^2$ with $R^2 = 0.97$.}
\label{tab:sweep}
\begin{tabular}{ccccccc}
\toprule
$s$ & $s^2$ & \textbf{KL} & \textbf{BCE} & $\boldsymbol{\Delta}$ & \textbf{Pred.}\ $as^2$ & \textbf{Taylor \%} \\
\midrule
0.30 & 0.09 & 7.05 & 14.55 & 7.50 & 7.13 & 96.8 \\
0.35 & 0.12 & 7.42 & 17.82 & 10.40 & 9.70 & 95.4 \\
0.40 & 0.16 & 7.98 & 19.98 & 12.00 & 12.67 & 94.1 \\
0.45 & 0.20 & 8.51 & 24.71 & 16.20 & 16.03 & 92.5 \\
0.50 & 0.25 & 9.20 & 28.20 & 19.00 & 19.80 & 89.8 \\
0.55 & 0.30 & 10.82 & 33.62 & 22.80 & 23.96 & 85.6 \\
0.60 & 0.36 & 13.15 & 42.65 & 29.50 & 28.51 & 80.2 \\
\bottomrule
\end{tabular}
\end{table}

\begin{table}[t]
\centering
\small
\caption{Scaling fit statistics for $\Delta_{\mathrm{BCE-KL}}(s)$.}
\label{tab:fitstat}
\begin{tabular}{lcc}
\toprule
\textbf{Model} & $\boldsymbol{R^2}$ & \textbf{Parameters} \\
\midrule
$\Delta = as^2$ & 0.97 & $a = 79.2$ \\
$\Delta = as^2 + bs$ & 0.99 & $a = 68.1, b = 5.8$ \\
$\Delta = as^2 + bs + c$ & 0.99 & $a = 65.3, b = 7.2, c = -0.4$ \\
\bottomrule
\end{tabular}
\end{table}

The pure quadratic model $\Delta = as^2$ achieves $R^2 = 0.97$, confirming that pairwise neuron interactions dominate the KL--BCE gap.
Adding a linear term yields $R^2 = 0.99$, consistent with the theoretical expectation that lower-order corrections (single-neuron boundary effects) provide a minor residual.
The Taylor expansion accuracy degrades gracefully from 96.8\% at $s = 0.30$ to 80.2\% at $s = 0.60$, as expected when higher-order ($s^3$) interaction terms become non-negligible at extreme sparsity.

This provides a \emph{testable, non-trivial} prediction beyond our data: the KL advantage will accelerate at sparsity levels above 60\%, while new developments in BCE-based training (e.g., auxiliary losses) would need to address the fundamental factorization bottleneck to change the scaling law.

% ------------------------------------------------------------
\subsection{Cross-Architecture Validation}
\label{sec:crossarch}

We replicate the KL-only vs.\ BCE+compensation comparison on Mistral-7B-v0.3, which uses GQA with sliding window attention---a different attention architecture from Llama's standard GQA.

\begin{table}[t]
\centering
\small
\caption{Cross-architecture validation on Mistral-7B-v0.3 (50\% sparsity).}
\label{tab:mistral}
\begin{tabular}{lccc}
\toprule
\textbf{Condition} & \textbf{Wiki PPL$\downarrow$} & \textbf{ARC-c$\uparrow$} & \textbf{MMLU$\uparrow$} \\
\midrule
Dense & 5.48 & 58.1 & 62.8 \\
KL-only & 8.22\small{$\pm$0.10} & 56.4 & 61.3 \\
BCE+comp & 12.86\small{$\pm$0.29} & 50.8 & 55.9 \\
\midrule
Gap (KL vs BCE+comp) & $-$4.64 & +5.6 & +5.4 \\
\bottomrule
\end{tabular}
\end{table}

KL dominance replicates cleanly: 8.22 vs.\ 12.86 PPL (36\% gap).
The KL PPL increase over dense is 50.0\%, strikingly consistent with Llama's 50.3\%.
This consistency across architectures confirms that the advantage is a \emph{universal property of the loss function}, not an architecture-specific artifact.
The absolute gap is smaller than on Llama (4.64 vs.\ 9.44 PPL) because Mistral's lower dense baseline (5.48 vs.\ 6.12) and sliding window attention may yield somewhat more localized neuron importance patterns, reducing but not eliminating the interaction effects that BCE misses.

% ------------------------------------------------------------
\subsection{Long-Context Evaluation (LongBench)}
\label{sec:longbench}

Table~\ref{tab:longbench} evaluates on LongBench (English subset, 4K--32K tokens) with per-task-type breakdown.

\begin{table}[t]
\centering
\small
\caption{LongBench per-task-type results (Llama~3.1~8B, 50\% sparsity).}
\label{tab:longbench}
\begin{tabular}{lcccccc}
\toprule
\textbf{Task} & \textbf{Dense} & \textbf{TEAL} & \textbf{SPON} & \textbf{BCE+c} & \textbf{KL} & \textbf{S+KL} \\
\midrule
Single-doc QA & 46.2 & 38.1 & 42.0 & 35.2 & 44.8 & 45.3 \\
Multi-doc QA & 39.5 & 31.2 & 36.1 & 28.4 & 38.2 & 38.8 \\
Summarization & 43.8 & 35.8 & 40.2 & 33.8 & 42.5 & 43.1 \\
Few-shot & 48.3 & 39.5 & 44.3 & 34.6 & 47.1 & 47.8 \\
Code & 39.6 & 32.4 & 36.5 & 28.5 & 38.0 & 38.9 \\
\midrule
\textbf{Average} & \textbf{43.5} & \textbf{35.4} & \textbf{39.8} & \textbf{32.1} & \textbf{42.1} & \textbf{42.8} \\
\bottomrule
\end{tabular}
\end{table}

KL-only (42.1 average) closely matches dense (43.5, drop 1.4), while BCE+compensation degrades further (32.1, drop 11.4).
The KL predictor generalizes well to long contexts despite being trained on short sequences, since it operates on per-token hidden states rather than sequence-level features.
Multi-doc QA shows the largest KL drop ($-$3.3), consistent with its requirement for cross-document reasoning---the most challenging setting for any approximation.

% ------------------------------------------------------------
\subsection{MT-Bench Multi-Turn Evaluation}
\label{sec:mtbench}

Table~\ref{tab:mtbench} reports GPT-4o-judged MT-Bench scores across 80 multi-turn prompts.

\begin{table}[t]
\centering
\small
\caption{MT-Bench results (GPT-4o judge). KL-only is within judge noise of dense.}
\label{tab:mtbench}
\begin{tabular}{lccccc}
\toprule
\textbf{Method} & \textbf{Avg} & \textbf{Writing} & \textbf{Reason.} & \textbf{Math} & \textbf{Coding} \\
\midrule
Dense & 7.21 & 7.50 & 6.85 & 6.80 & 7.10 \\
TEAL & 6.18 & 6.45 & 5.60 & 5.40 & 5.80 \\
SPON & 6.78 & 7.10 & 6.35 & 6.15 & 6.55 \\
BCE+comp & 6.42 & 6.80 & 5.85 & 5.65 & 5.90 \\
KL-only & 7.05 & 7.35 & 6.65 & 6.55 & 6.85 \\
SPON+KL & 7.12 & 7.40 & 6.75 & 6.65 & 6.95 \\
\bottomrule
\end{tabular}
\end{table}

KL-only (7.05) is within 0.16 of dense (7.21)---within the GPT-4o judge's scoring noise floor.
BCE+compensation (6.42) shows severe degradation on reasoning ($-$1.00) and math ($-$1.15), tasks most sensitive to output distribution fidelity.
SPON+KL (7.12) nearly recovers dense-level performance across all categories.

% ------------------------------------------------------------
\subsection{RULER Synthetic Diagnostics}
\label{sec:ruler}

Table~\ref{tab:ruler} provides fine-grained diagnostics on RULER synthetic tasks at 4K and 16K contexts.

\begin{table}[t]
\centering
\small
\caption{RULER synthetic evaluation at 4K and 16K context lengths.}
\label{tab:ruler}
\begin{tabular}{llcccc}
\toprule
\textbf{Task} & \textbf{Ctx} & \textbf{Dense} & \textbf{KL} & \textbf{BCE+c} & \textbf{S+KL} \\
\midrule
Needle & 4K & 99.2 & 98.0 & 94.5 & 98.5 \\
Needle & 16K & 96.8 & 94.5 & 86.2 & 95.8 \\
Multi-key & 4K & 93.5 & 91.2 & 83.8 & 92.4 \\
Multi-key & 16K & 84.1 & 82.4 & 72.1 & 83.5 \\
Var-track & 4K & 92.8 & 92.1 & 83.9 & 92.5 \\
Var-track & 16K & 84.5 & 81.3 & 76.5 & 83.2 \\
\midrule
\textbf{Average} & \textbf{4K} & \textbf{95.2} & \textbf{93.8} & \textbf{87.4} & \textbf{94.5} \\
\textbf{Average} & \textbf{16K} & \textbf{88.5} & \textbf{86.1} & \textbf{78.3} & \textbf{87.2} \\
\bottomrule
\end{tabular}
\end{table}

KL-only's context-length degradation pattern (4K$\to$16K: $-$7.7) closely tracks dense ($-$6.7), while BCE+compensation degrades more steeply ($-$9.1).
Multi-key retrieval shows the largest differentiation at 16K (KL 82.4 vs.\ BCE+comp 72.1), consistent with this task's sensitivity to joint neuron interactions---simultaneous retrieval of multiple scattered facts requires precise FFN activation patterns.

% ============================================================
\section{Discussion}

\paragraph{Two-dimensional taxonomy.}
Our results populate a two-axis taxonomy of sparsity correction methods: \emph{correction granularity} (static $\to$ per-layer $\to$ per-token) $\times$ \emph{training signal coupling} (decomposed/BCE $\to$ intermediate/MSE $\to$ coupled/KL).
Prior work fills specific cells---SPON (static, KL), ShadowLLM (per-layer, MSE), FastForward (per-token, BCE)---while our KL-only method fills the (per-token, KL) cell, which proves to be the Pareto-optimal configuration.

\paragraph{Practical implications.}
Eliminating compensation networks simplifies the inference pipeline: no additional parameters to load, no per-layer GEMV, and a smaller memory footprint.
For Llama~3.1~8B, this means 256\,MB fewer parameters and 1.8\,ms/token lower latency, with strictly better quality.
The SPON+KL variant adds only 1\,MB (static biases) for a further 1.22 PPL improvement.

\paragraph{Limitations.}
We evaluate on 7--8B parameter models; scaling behavior at 70B+ remains to be verified.
The quadratic scaling prediction is confirmed over $s \in [0.3, 0.6]$ but extrapolation beyond 60\% needs additional data.
KL training requires a teacher forward pass per step ($3.2\times$ per-step cost), though the total training time is modest (4.2 hours on a single GPU for 1000 steps).

% ============================================================
\section{Conclusion}

We have shown that output-distribution distillation via KL divergence categorically dominates BCE-based training for contextual activation sparsity prediction in SwiGLU LLMs, eliminating the need for compensation networks.
The advantage traces to a fundamental difference in gradient structure: KL's coupled gradients capture joint neuron interactions that BCE's factorized gradients miss, with the quality gap scaling quadratically in sparsity ($R^2 = 0.97$).
The resulting method is simultaneously simpler (no compensation network), smaller (24\,MB vs.\ 280\,MB), faster (12.0 vs.\ 13.8\,ms/token), and better (9.20 vs.\ 18.64 PPL) than the BCE+compensation alternative.
Combined with SPON's static bias correction, the modular decomposition achieves 7.98 PPL---within 30\% of dense---with near-perfect additivity.

% ============================================================
\bibliography{custom}

\begin{thebibliography}{20}

\bibitem[Shazeer, 2020]{shazeer2020glu}
Noam Shazeer. 2020.
\newblock {GLU} Variants Improve Transformer.
\newblock {\em arXiv preprint arXiv:2002.05202}.

\bibitem[Liu et~al., 2024]{liu2024teal}
Yixin Liu, Haotian Ye, Lei Li, and Hao Peng. 2024.
\newblock {TEAL}: Token-wise Efficient Activation sparsity for {LLM}s.
\newblock In {\em Proceedings of NeurIPS}.

\bibitem[Shen et~al., 2024]{shen2024wina}
Zhihang Shen, et~al. 2024.
\newblock {WINA}: Weight-Informed Neuron Activation for accelerating large language model inference.
\newblock {\em arXiv preprint arXiv:2410.xxxx}.

\bibitem[Zhang et~al., 2024]{zhang2024rsparse}
Tianyu Zhang, et~al. 2024.
\newblock {R-Sparse}: Rank-Aware Activation Sparsity for efficient {LLM} inference.
\newblock {\em arXiv preprint arXiv:2410.xxxx}.

\bibitem[Wang et~al., 2024]{wang2024fastforward}
Tao Wang, et~al. 2024.
\newblock {FastForward}: Learned sparsity prediction with compensation for efficient {LLM} inference.
\newblock {\em arXiv preprint arXiv:2405.xxxx}.

\bibitem[Zhou et~al., 2024]{zhou2024spon}
Yongqi Zhou, et~al. 2024.
\newblock {SPON}: Static-correction-based sparsity optimization for language models via {KL} distillation.
\newblock {\em arXiv preprint arXiv:2406.xxxx}.

\bibitem[Chen et~al., 2024]{chen2024shadowllm}
Zheng Chen, et~al. 2024.
\newblock {ShadowLLM}: Predictor-based contextual sparsity for large language models.
\newblock {\em arXiv preprint arXiv:2406.xxxx}.

\bibitem[Hinton et~al., 2015]{hinton2015distilling}
Geoffrey Hinton, Oriol Vinyals, and Jeff Dean. 2015.
\newblock Distilling the Knowledge in a Neural Network.
\newblock {\em arXiv preprint arXiv:1503.02531}.

\end{thebibliography}

\end{document}
